% ============================================================
% Section 4: Indian Mathematics (Slides 25-31)
% ~500 BCE -- 1200 CE
% ============================================================
\section{Indian Mathematics}

% ----------------------------------------------------------
% Slide 25: Section Divider
% ----------------------------------------------------------
{
\setbeamercolor{background canvas}{bg=indianSaffron}
\begin{frame}[plain,shrink=23]
  \centering
  \vspace{1cm}
  \textcolor{white}{\Huge\bfseries India}\\[0.5em]
  \textcolor{white!90}{\Large The Invention of Nothing (and Everything)}\\[0.8em]
  \textcolor{white!75}{\large $\sim$500 BCE -- 1200 CE}\\[1.5em]

  % Stylized zero with Brahmi numerals
  \visualplaceholder[5cm]{The number that changed everything}

  \vspace{0.8em}
  % Timeline marker
  \visualplaceholder[5cm]{Diagram}

  \note{
    Section transition. India's contributions to mathematics are often under-appreciated
    in Western curricula. The decimal system, zero, negative numbers, trigonometry,
    infinite series --- all originate here.
  }
\end{frame}
}

% ----------------------------------------------------------
% Slide 26: Early Indian Mathematics
% ----------------------------------------------------------
\begin{frame}[shrink=9]{Early Indian Mathematics --- The Sulba Sutras}

  \begin{columns}[T]
    \begin{column}{0.52\textwidth}
      \begin{itemize}
        \item \textbf{Sulba Sutras} (c.\ 800--200 BCE)
        \item Geometric rules for Vedic \alert{altar construction}
        \item Include \textbf{Pythagorean triples} --- independently of Babylon and Greece
        \item Approximation of $\sqrt{2}$:
          {\small
          \[
            \sqrt{2} \approx 1 + \tfrac{1}{3} + \tfrac{1}{3 \cdot 4} - \tfrac{1}{3 \cdot 4 \cdot 34}
            \approx 1.4142156\ldots
          \]}
          (true value: $1.4142135\ldots$)
      \end{itemize}

      \vspace{0.3em}
      \begin{insightbox}
        Pythagorean triples appear in at least THREE independent traditions:
        Babylonian, Indian, and Greek.
      \end{insightbox}
    \end{column}

    \begin{column}{0.45\textwidth}
      % Altar construction geometry
      \visualplaceholder[5cm]{Falcon-shaped altar (simplified)}

      \vspace{0.3em}
      \visualplaceholder[4.5cm]{Diagram from the Sulba Sutras showing altar construction with geometric proportions}
    \end{column}
  \end{columns}

  \note{
    The Sulba Sutras are religious texts (rules for constructing fire altars),
    but they contain sophisticated geometry. The 3-4-5 triple appears here too.
    \verifytag{Dating of Sulba Sutras varies. The Baudhayana Sulba Sutra is
    generally dated c.\ 800 BCE but some scholars place it later.}

    NARRATIVE BEAT -- THE PATTERN: ``This is the third time we have seen this.
    Pythagorean triples in Babylon (Slide 7), in Greece (Pythagoras), and now
    in India -- independently. People who had no contact with each other,
    discovering the same truth. What does that tell us about mathematics?
    It tells us that these patterns are OUT THERE, waiting to be found.
    Mathematics is not invented -- it is discovered.'' [Pause for effect.]
  }
\end{frame}

% ----------------------------------------------------------
% Slide 27: Aryabhata
% ----------------------------------------------------------
\begin{frame}[shrink=16]{Aryabhata}

  \begin{columns}[T]
    \begin{column}{0.55\textwidth}
      \begin{personbox}[Aryabhata (Aryabhata I)]{indianSaffron}
        \small
        \textbf{Dates:} 476 -- 550 CE\\
        \textbf{Origin:} Kusumapura (likely modern Patna), Gupta Empire, India\\
        \textbf{Contribution:} Wrote \emph{Aryabhatiya} (499 CE) --- trigonometric tables,
        place-value system, approximation of $\pi$, Earth's rotation
      \end{personbox}

      \vspace{0.4em}
      \begin{itemize}
        \item $\pi \approx 3.1416$ --- remarkably accurate
        \item Used a \alert{place-value number system}
        \item Tabulated \textbf{sine values} (the first systematic trigonometric table)
        \item Proposed that the \textbf{Earth rotates on its axis}
        \item \textbf{Kuttaka method} for indeterminate equations
      \end{itemize}
    \end{column}

    \begin{column}{0.42\textwidth}
      \visualplaceholder[4.5cm]{Statue of Aryabhata (IUCAA Pune) or page from a manuscript of \emph{Aryabhatiya} (Wikimedia Commons)}

      \vspace{0.5em}
      % Sine table illustration
      \visualplaceholder[5cm]{Data points}
    \end{column}
  \end{columns}

  \note{
    Aryabhata is one of the most important mathematicians in history.
    His place-value system is the ancestor of the numerals we use today.
    His value of $\pi$ was not bettered in India for centuries.
    \verifytag{Location of statue: IUCAA Pune --- confirm.}
  }
\end{frame}

% ----------------------------------------------------------
% Slide 28: Brahmagupta
% ----------------------------------------------------------
\begin{frame}[shrink=12]{Brahmagupta --- The Rules for Zero}

  \begin{columns}[T]
    \begin{column}{0.52\textwidth}
      \begin{personbox}[Brahmagupta]{indianSaffron}
        \small
        \textbf{Dates:} 598 -- c.\ 668 CE\\
        \textbf{Origin:} Bhinmal, Rajasthan, India\\
        \textbf{Contribution:} \emph{Brahmasphutasiddhanta} (628 CE) --- first to
        formalize rules for zero and negative numbers; general quadratic formula;
        cyclic quadrilateral formula
      \end{personbox}

      \vspace{0.3em}
      \begin{itemize}
        \item \alert{First rules for arithmetic with zero:}\\
          $a + 0 = a$, \quad $a \times 0 = 0$, \quad $0 + 0 = 0$
        \item Negative $\times$ negative $=$ \textbf{positive}
        \item General solution to the \textbf{quadratic equation}
        \item Said $0 \div 0 = 0$ \textrightarrow\ mathematically wrong, but historically important
      \end{itemize}
    \end{column}

    \begin{column}{0.45\textwidth}
      % Rules for zero - visual
      \visualplaceholder[5cm]{Rules in boxes}

      \vspace{0.3em}
      % Cyclic quadrilateral
      \visualplaceholder[5cm]{Cyclic quadrilateral}
    \end{column}
  \end{columns}

  \note{
    ``Brahmagupta wrote the rules for zero that you use every day --- thirteen centuries ago.''
    \verifytag{Brahmagupta's treatment of $0/0$ (he said $0/0 = 0$) is mathematically
    incorrect but historically important. Mention this nuance.}
    These rules would travel to the Islamic world via al-Khwarizmi and then to Europe.
  }
\end{frame}

% ----------------------------------------------------------
% Slide 29: Bhaskara II
% ----------------------------------------------------------
\begin{frame}[shrink=21]{Bhaskara II (Bhaskaracharya)}

  \begin{columns}[T]
    \begin{column}{0.55\textwidth}
      \begin{personbox}[Bhaskara II --- ``Bhaskara the Teacher'']{indianSaffron}
        \small
        \textbf{Dates:} 1114 -- 1185 CE\\
        \textbf{Origin:} Bijjada Bida (likely modern Karnataka), India\\
        \textbf{Contribution:} Wrote \emph{Lilavati} (mathematics) and
        \emph{Bijaganita} (algebra); solved Pell's equation before Pell;
        preliminary concepts related to infinitesimals
      \end{personbox}

      \vspace{0.3em}
      \begin{itemize}
        \item \emph{Lilavati} --- a math textbook \alert{addressed to his daughter}
        \item One of the earliest textbooks written for a specific student
        \item Solved $x^2 - ny^2 = 1$ (``Pell's equation'') \textbf{before} Pell
        \item Preliminary ideas related to \textbf{derivatives} and infinitesimals
      \end{itemize}
    \end{column}

    \begin{column}{0.42\textwidth}
      \visualplaceholder[4.5cm]{Opening verse of \emph{Lilavati} in Devanagari with English translation}

      \vspace{0.3em}
      % Pell's equation example
      \visualplaceholder[5cm]{Diagram}
    \end{column}
  \end{columns}

  \note{
    ``He addressed his math textbook to his daughter Lilavati --- making it one
    of the earliest math textbooks written with a specific student in mind.''
    \verifytag{The story of Lilavati being his daughter and the astrological
    tale around her is traditional but may be legend.}
    The Pell equation example is stunning: a 12th-century Indian mathematician
    found solutions with 10-digit numbers.
  }
\end{frame}

% ----------------------------------------------------------
% Slide 30: The Kerala School
% ----------------------------------------------------------
\begin{frame}[shrink=25]{The Kerala School --- Calculus Before Calculus?}

  \begin{columns}[T]
    \begin{column}{0.52\textwidth}
      \begin{personbox}[Madhava of Sangamagrama]{indianSaffron}
        \small
        \textbf{Dates:} c.\ 1340 -- c.\ 1425 CE\\
        \textbf{Origin:} Sangamagrama (modern Irinjalakuda), Kerala, India\\
        \textbf{Contribution:} Founder of the Kerala school; infinite series for
        $\pi$, sine, and cosine --- BEFORE European discovery
      \end{personbox}

      \vspace{0.3em}
      Also: \textbf{Nilakantha Somayaji} (1444--1544 CE, Kerala) ---
      wrote \emph{Tantrasamgraha} (1501), extended Madhava's work.

      \vspace{0.3em}
      \begin{itemize}
        \item Infinite series for $\pi$ (Madhava--Leibniz):
          \[
            \frac{\pi}{4} = 1 - \frac{1}{3} + \frac{1}{5} - \frac{1}{7} + \cdots
          \]
        \item Power series for $\sin x$ and $\cos x$
        \item \textbf{200+ years} before Newton and Leibniz
      \end{itemize}
    \end{column}

    \begin{column}{0.45\textwidth}
      % Timeline comparison
      \visualplaceholder[5cm]{Kerala school}

      \vspace{0.3em}
      \begin{insightbox}
        Was this transmitted to Europe or independently rediscovered? An open question in the history of mathematics.
      \end{insightbox}
    \end{column}
  \end{columns}

  \note{
    This is one of the most striking priority stories in all of mathematics.
    \verifytag{The extent to which Kerala school results were transmitted to Europe
    is debated. Most historians believe independent rediscovery, but some argue
    for possible transmission via Jesuit missionaries. Present as an open question.}
    \verifytag{Nilakantha's death year --- 1544 CE is the standard but verify.}
  }
\end{frame}

% ----------------------------------------------------------
% Slide 31: Indian Mathematics Legacy
% ----------------------------------------------------------
\begin{frame}{Indian Mathematics --- A Legacy for the World}

  \begin{columns}[T]
    \begin{column}{0.45\textwidth}
      \textbf{India gave the world:}
      \begin{itemize}
        \item The \alert{decimal place-value system}
        \item \textbf{Zero} as a number
        \item \textbf{Negative numbers} with formal rules
        \item Trigonometric functions (\textbf{sine})
        \item \textbf{Infinite series}
        \item Solutions to Pell's equation
        \item The foundations of the number system we use today
      \end{itemize}
    \end{column}

    \begin{column}{0.52\textwidth}
      % Flow diagram: Indian numerals -> Arabic -> World
      \visualplaceholder[5cm]{Indian numerals}

      \vspace{0.5em}
      {\small The numeral system you use every day\\
      was invented in India, transmitted through\\
      the Islamic world, and adopted by Europe.}
    \end{column}
  \end{columns}

  \vspace{0.3em}
  \begin{center}
    \textcolor{chineseRed}{\emph{``Meanwhile, on the other side of Asia, a parallel mathematical
    tradition was flourishing\ldots''}}
  \end{center}

  \note{
    Summary and transition. The key takeaway is the FLOW of knowledge:
    India $\to$ Islamic world $\to$ Europe.
    ``Indian mathematical ideas would soon be taken up, extended, and transformed
    by scholars writing in Arabic\ldots'' But first, we look east to China.
  }
\end{frame}
